\subsection{对合 Banach 代数的包络星代数}

引理 16. --- 设 A 是一个对合代数，p 是 A 上的一个半范数。下列条件等价：

\begin{enumerate}
\def\labelenumi{(\roman{enumi})}
\item
  我们有 \(p(xy) \leqslant p(x)p(y)\)，\(p(x^{*}) = p(x)\) 且 \(p(x)^{2} = p(x^{*}x)\) 对一切 \(x, y \in \mathrm{A}\) 成立；
\item
  A 中满足 \(p(x) = 0\) 的元素 x 组成的集合 R 是 A 的一个自伴双边理想，且由 p 在 A/R 上诱导的半范数使 A/R 成为一个对合赋范代数，其完备化是一个星代数（C*-代数）；
\item
  存在一个星代数 B 及从 A 到 B 的一个对合代数态射 \(\varphi\)，使得 \(p(x) = \|\varphi(x)\|\) 对一切 \(x\in\mathbf{A}\) 成立。
\end{enumerate}

蕴含关系 (i) \(\Longrightarrow\) (ii) \(\Longrightarrow\) (iii) \(\Longrightarrow\) (i) 都是初等的。

满足引理 16 中条件的半范数将称为对合代数 A 上的星半范数。

设 A 是一个赋范对合代数，S 是 A 上星半范数组成的集合。我们有 \(p(x) \leqslant \|x\|\) 对一切 \(x \in A\) 及每个 \(p \in S\) 成立 (I, p.~104, 命题 2)。映射 \(x \mapsto \|x\|_* = \sup_{p \in S} p(x)\) 是 A 上的一个星半范数。它是 A 上最大的星半范数。

设 R 是由 \(x \in A\) 中满足 \(\|x\|_{*} = 0\) 的元素组成的集合。它是 A 的一个闭双边理想。我们用 Stell(A) 表示 A/R 对于由 \(x \mapsto \|x\|_{*}\) 诱导的范数的完备化星代数（引理 16, (ii)）。从 A 到 Stell(A) 中的典范映射是连续的，它的像在 Stell(A) 中稠密，且它的核等于 R。

定义 8. --- 星代数 Stell(A) 称为赋范对合代数 A 的包络星代数。

若 A 是交换的，则 Stell(A) 是交换的；若 A 是幺的，则 Stell(A) 是幺的。

命题 20. --- 设 A 是一个赋范对合代数，j 是从 A 到 Stell(A) 中的典范态射。对每个星代数 B 及每个从 A 到 B 的对合代数态射 \(\varphi\)，存在唯一的从 Stell(A) 到 B 中的星代数态射 \(\varphi'\) 使 \(\varphi = \varphi' \circ j\)。

我们用 \(x \mapsto \|x\|_*\) 表示 Stell(A) 上的范数。设 R 是 \(j\) 的核。映射 \(x \mapsto \| \varphi(x) \|\) 是 A 上的一个星半范数。于是 \(\| \varphi(x) \| \leqslant \| x \|_*\) 对一切 \(x \in A\) 成立。因而态射 \(\varphi\) 通过取商定义了从 A/R 到 B 的一个连续态射，它由连续性拓展为从 Stell(A) 到 B 的一个态射 \(\varphi'\)，满足 \(\varphi = \varphi' \circ j\)。\(\varphi'\) 的唯一性源于 \(j\) 的像在 Stell(A) 中稠密这一事实。

推论. --- 设 A 是一个交换对合 Banach 代数，j 是从 A 到 Stell(A) 中的典范态射。映射 \(\mathsf{X}(j)\) 是从 \(\mathsf{X}(\mathrm{Stell}(\mathrm{A}))\) 到子空间 \(\mathsf{X}(\mathrm{A})_{h}\)（由 A 的 Hermite 特征标组成的 \(\mathsf{X}(\mathrm{A})\) 的子空间）上的一个同胚。

A 的 Hermite 特征标就是从 A 到星代数 C 的对合代数态射。于是命题 20 蕴含 \(\mathsf{X}(j)\) 是从 \(\mathsf{X}(\mathrm{Stell}(\mathsf{A}))\) 到 \(\mathsf{X}(\mathsf{A})_h\) 上的一个双射。由于 \(\mathsf{X}(j)\) 是到其像上的一个同胚（参见 I, p.~10），推论得证。

我们把 \(\mathsf{X}(\mathrm{Stell}(\mathbf{A}))\) 与 \(\mathsf{X}(\mathbf{A})_h\) 通过映射 \(\mathsf{X}(j)\) 恒等同。对每个 \(x\in \mathbf{A}\)，映射 \(\mathcal{G}_{\mathrm{Stell(A)}}(j(x))\) 在 \(\mathsf{X}(\mathbf{A})_h\) 上不是别的，正是 \(\mathcal{G}_{\mathbf{A}}(x)\)。

命题 21. --- 设 A 是一个对合 Banach 代数，j 是从 A 到 Stell(A) 中的典范态射。A 的根式包含在 j 的核中。

设 \(x\) 是 A 的根式的一个元素。那么 \(x^{*}x\) 属于 A 的根式，因此 \(\mathrm{Sp}_{\mathrm{A}}^{\prime}(x^{*}x) = \{0\}\) (I, p.~5, 注记 3)。由于 \(\mathrm{Sp}_{\mathrm{Stell(A)}}^{\prime}(j(x^{*}x)) \subset \mathrm{Sp}_{\mathrm{A}}^{\prime}(x^{*}x)\)，我们于是有 \(\mathrm{Sp}_{\mathrm{Stell(A)}}^{\prime}(j(x)^{*}j(x)) = \{0\}\)，从而 \(\| j(x)\| ^2 = \| j(x)^{*}j(x)\| = \varrho (j(x)^{*}j(x)) = 0\) (I, p.~104, 公式 (2))，因此 \(j(x) = 0\)。
% semantic-fragments: legacy-input-seam
\relax{}
